\section{Models and the shared FiLM intervention}
\label{sec:models}
The paper contains two linked studies. The \emph{benchmark study} compares representative predictors
under identical COAST inputs. The \emph{conditioning study} changes one factor inside each predictor:
whether global or local morphology is injected through FiLM. Keeping these studies separate prevents
a strong backbone from being mistaken for evidence that conditioning itself works.

\subsection{Give every method the same inputs}
Every method receives identical frozen 1024-dimensional UNI patch features, coordinates, fold-specific
panels, and normalized targets. The main paper evaluates linear ridge and random-forest probes;
ST-Net, DeepSpaCE, an MLP probe, HisToGene, Hist2ST, a Gene-DML-style multi-level predictor,
and the whole-slide STFlow generator. This covers
spot-independent regression, slide transformers, graph/spatial models, and generative prediction.
The appendix contains the complete screening table, including retrieval and multi-resolution models.
Training budgets and validation rules are matched where model
interfaces permit; all deviations are disclosed rather than silently harmonized.

\subsection{Change one factor: morphology conditioning}
Organ labels are unavailable for a novel organ as learned categories. We therefore test whether
continuous histology, which remains available at inference, is a better adaptation signal. We insert
the same zero-initialized FiLM principle into ST-Net, DeepSpaCE, HisToGene, Hist2ST, an MLP probe,
and the Gene-DML-style predictor,
comparing each model only with its own matched unconditioned version. The
\emph{global} intervention uses the masked slide mean; the \emph{local} intervention also uses the
mean feature in each spot's spatial neighborhood. Results for additional architectures are retained
in the appendix to make the full screening record visible.

Before comparing results, we classify each backbone by the context already used in its original
prediction path (Table~\ref{tab:context}). This supports a predeclared question: does FiLM help most
when a backbone does not already aggregate both neighborhood and slide context? The taxonomy is based
on architecture, not observed performance.

\begin{table}[t]
\centering\small
\setlength{\tabcolsep}{3.5pt}
\resizebox{\columnwidth}{!}{%
\begin{tabular}{lccc}
\toprule
Backbone & Spot & Neighborhood & Slide/global\\
\midrule
ST-Net / DeepSpaCE / MLP & yes & no & no\\
HisToGene & yes & limited & yes\\
Hist2ST & yes & yes & limited\\
Gene-DML-style & yes & yes & yes\\
STFlow & yes & yes & input only\\
\bottomrule
\end{tabular}
}
\caption{Context used by each original prediction path. The taxonomy is defined independently of
the FiLM results.}
\label{tab:context}
\end{table}

For the geometric STFlow reference, an adaLN-Zero head in each transformer block maps an invariant conditioner $\mathbf c$
to shift, scale, and residual gate parameters:
\begin{equation}
\tilde{\mathbf h}=\mathbf h\odot(1+\gamma(\mathbf c))+\beta(\mathbf c),\qquad
\mathbf h' = \mathbf h+a(\mathbf c)\odot F(\tilde{\mathbf h}).
\end{equation}
The heads are zero-initialized. The global variant uses the masked mean of slide patch features. The
local variant adds, for each spot, the mean feature among its spatial $k$ nearest neighbors. Both
conditioners are invariant to coordinate rotation and translation. Modulation acts only on STFlow's
invariant scalar stream and gates the output of geometric attention; it never modifies frame or edge
geometry. A numerical transformation test verifies prediction differences below $10^{-6}$.

For every backbone, the controlled comparison is conditioned versus unconditioned training with identical
features, folds, objectives, seeds, and optimization budgets. Parameter-matched unconditioned controls
separate conditioning information from added capacity. The larger-capacity STFlow model from
exploratory experiments is not a headline method and is reported only as a secondary scaling analysis.
